\chapter{기계는 밤에 무엇을 하는가}
% full version: chapters/20_sleep.tex

잠은 기계를 끄는 것이 아니라 모드를 바꾸는 것이다. 방송국을 보면 알 수 있다. 깊은 잠에서는 아세틸콜린, 노르아드레날린, 세로토닌 등 거의 모든 방송국의 소리가 줄어든다. 꿈을 꾸는 잠에서는 아세틸콜린이 낮 수준으로 돌아오고, 노르아드레날린과 세로토닌은 거의 침묵하며, 근육으로 가는 출력은 브레이크로 꺼진다. 그래서 꿈속에서 무언가에게 쫓길 때 우리 몸은 실제로 뛰지 않는다. 이것은 모두 측정된 사실이다.

\section*{잘 시간은 언제인가}

깨어 있는 동안 하루 종일 충전되고 밤에 방전되는 커패시터를 떠올려 보자. 전하가 곧 “수면 압력”이다. 전하가 위쪽 문턱값에 이르면 잠들고, 아래쪽 문턱값까지 내려가면 깬다. 그리고 두 문턱값은 앞 장의 하루 리듬을 따라 부드럽게 오르내린다. 이 단순한 회로는 저절로 16시간 깨어 있고 8시간 자는 주기에 이른다. 또 밤을 새운 뒤에 두 배로 오래 자는 것이 아니라 더 깊이 자는 까닭도 설명한다. 가득 찬 커패시터는 더 빨리 방전된다. “전하”의 후보는 뇌가 일하는 동안 쌓이는 물질이다. 그것이 정말 이 물질인지는 아직 가설이다.

\pencilfig{120}{%
% figures/sleep_{S,H,L}.dat from models/sleep_models.py, t in hours
\begin{scope}[x=0.1cm, y=2.6cm]
\foreach \a/\b in {0/4.3,21.2/29.3,45.4/53.4,69.5/72}{\fill[pcBlue, opacity=0.1] (\a,0) rectangle (\b,1);}
\draw[pen] (0,0) -- (72,0);
\draw[penlite=pcRed, densely dashed] plot file {figures/sleep_H.dat};
\draw[penlite=pcGreen, densely dashed] plot file {figures/sleep_L.dat};
\draw[pcBlue, line width=1.0pt] plot file {figures/sleep_S.dat};
\foreach \t/\l in {0/0,24/1일,48/2일,72/3일}{\draw[pcGray, line width=0.5pt] (\t,0) -- (\t,-0.04); \node[lbl, anchor=base] at (\t,-0.16) {\l};}
\node[lbl, text=pcBlue] at (25.25,0.93) {수면}; \node[lbl, text=pcBlue] at (49.4,0.93) {수면};
\end{scope}
\node[lbl, text=pcRed, anchor=west] at (7.3,2.0) {잠들기};
\node[lbl, text=pcGreen!70!black, anchor=west] at (7.3,0.62) {깨기};
\node[lbl, text=pcBlue, anchor=west] at (7.3,1.35) {수면 압력};
}{수면 압력은 낮에 쌓이고 밤에 녹아내린다. 위쪽 문턱값에서 잠들고 아래쪽 문턱값에서 깨며, 두 문턱값은 하루 주기를 따라 흔들린다.}

\section*{느린 그네}

깊은 잠에서는 뇌의 넓은 영역 전체가 한꺼번에 켜졌다 꺼졌다를 되풀이한다. 1초에 한 번보다 드물다. 이미 익숙한 회로다. 스스로를 흥분시키다가 지치는 무리, 곧 발진기다. 낮에는 아세틸콜린이 뉴런의 피로를 억눌러서 같은 무리가 고르게 일한다. 밤에는 아세틸콜린이 적어서 네트워크가 흔들리기 시작한다.

\pencilfig{20}{\pencilart{20}}{밤에는 네트워크가 느리게 흔들린다. 거의 모든 뉴런이 일하다가, 거의 모두가 침묵한다.}

\section*{빨리 감기}

가장 흥미로운 일은 이 그네 안에서 일어난다. 낮에 일어났던 뉴런 발화의 사슬이 다시 “재생”되는데, 몇 배나 빠르다. 기억 영역에서는 약 20배, 대뇌 피질에서는 몇 배다. 이 재생과 느린 그네의 맞물림을 인위적으로 강화하면 기억이 좋아진다. 이것은 인과를 확인한 실험이다. 왜 이런 일이 필요할까? 그럴듯한 가설은 이렇다. 느린 장기 기억은 반복된 재생을 옛것과 뒤섞어 배우며, 그래야 새것이 옛것을 지워 버리지 않는다. 엔지니어들은 이 문제를 안다. 새 과제로 학습시킨 신경망은 옛 예제를 되풀이해 보여 주지 않으면 옛 과제를 잊어버린다.

\section*{연결 다듬기}

측정된 것이 하나 더 있다. 잠을 자고 나면 뉴런 사이의 접점이 평균적으로 몇 퍼센트에서 수십 퍼센트까지, 각자의 크기에 비례해 작아진다. 모든 음량을 한 번에 조금씩 줄인 것과 같다. 비율은 그대로이고, 배운 것은 남으며, 내일 학습할 자리가 생긴다. 이것이 잠의 주된 임무라는 것은 가설이다.

\section*{꿈과 청소}

꿈을 꾸는 잠은 벤치 테스트와 비슷하다. 기계는 최대 출력으로 돌아가지만, 입력은 내부 신호로 바뀌고 출력은 끊겨 있다. 이것이 왜 필요한지는 모른다. 잠자는 동안 뇌가 노폐물을 “씻어 낸다”는 생각은 아직 논란거리다. 찬성하는 실험도, 반대하는 실험도 있다.

가장 놀라운 점은 잠이 국소적일 수 있다는 것이다. 오래 깨어 있으면 깨어 있는 사람의 대뇌 피질 일부가 한순간 “잠들어”, 일하는 도중에 조용해진다.

\fullversion{20}
